% ============================================================

% APC2 - Cálculo Numérico Aplicado
% Resolução analítica baseada na organização conceitual do APC1
% ============================================================
\documentclass[12pt,a4paper]{article}

\usepackage[utf8]{inputenc}
\usepackage[T1]{fontenc}
\usepackage[brazil]{babel}
\usepackage{amsmath,amssymb,mathtools}
\usepackage{geometry}
\usepackage{siunitx}
\usepackage{booktabs}
\usepackage{array}
\usepackage{enumitem}
\usepackage{float}
\usepackage{microtype}

\geometry{margin=2.5cm}
\sisetup{output-decimal-marker={,}}

\title{\textbf{Atividade para Casa \#2 (APC2)}\\
Cálculo Numérico Aplicado}
\author{}
\date{2026}

\begin{document}

\maketitle

\section*{Introdução}

A presente atividade tem como objetivo desenvolver a formulação analítica de um
sistema massa--mola--amortecedor com múltiplos graus de liberdade e estabelecer
a conexão entre o modelo físico e o polinômio característico que será utilizado
posteriormente no Programa para Casa \#2 (PPC2).

A organização da resolução segue a mesma lógica utilizada na APC1: parte-se da
equação governante, propõe-se uma forma de solução, substitui-se essa solução na
equação original, reorganizam-se os termos e, por fim, obtém-se a forma
matemática que será utilizada na etapa computacional.

A APC2 solicita, em particular, a obtenção do polinômio característico de ordem
$2N$ associado ao sistema e sua organização na forma padrão para implementação
computacional.

\section{Formulação geral do sistema}

Para um sistema com $N$ graus de liberdade, a equação do movimento fornecida no
enunciado é

\begin{equation}
\mathbf{M}\frac{d^2\mathbf{x}}{dt^2}
+
\mathbf{C}\frac{d\mathbf{x}}{dt}
+
\mathbf{K}\mathbf{x}
= \mathbf{0},
\label{eq:movimento}
\end{equation}

em que:

\begin{itemize}
    \item $\mathbf{M}$ é a matriz de massa;
    \item $\mathbf{C}$ é a matriz de amortecimento;
    \item $\mathbf{K}$ é a matriz de rigidez;
    \item $\mathbf{x}(t)$ é o vetor de deslocamentos.
\end{itemize}

Para $N$ graus de liberdade,

\begin{equation}
\mathbf{x}(t)=
\begin{bmatrix}
x_1(t) & x_2(t) & \cdots & x_N(t)
\end{bmatrix}^{T}.
\end{equation}

As matrizes $\mathbf{M}$, $\mathbf{C}$ e $\mathbf{K}$ possuem dimensão
$N\times N$.

\section{Substituição da solução proposta}

Assim como na APC1, adotamos uma solução exponencial. Para o problema com
múltiplos graus de liberdade, o enunciado propõe

\begin{equation}
\mathbf{x}(t)=\boldsymbol{\phi}e^{\lambda t},
\label{eq:solucao}
\end{equation}

em que $\boldsymbol{\phi}$ é um vetor constante associado à forma modal e
$\lambda$ é um parâmetro a ser determinado.

Derivando a Equação~\eqref{eq:solucao} uma vez em relação ao tempo,

\begin{equation}
\frac{d\mathbf{x}}{dt}
=
\lambda\boldsymbol{\phi}e^{\lambda t}.
\label{eq:primeira}
\end{equation}

Derivando novamente,

\begin{equation}
\frac{d^2\mathbf{x}}{dt^2}
=
\lambda^2\boldsymbol{\phi}e^{\lambda t}.
\label{eq:segunda}
\end{equation}

Substituindo as Equações~\eqref{eq:primeira} e
\eqref{eq:segunda} na Equação~\eqref{eq:movimento}, obtemos

\begin{equation}
\mathbf{M}
\left(\lambda^2\boldsymbol{\phi}e^{\lambda t}\right)
+
\mathbf{C}
\left(\lambda\boldsymbol{\phi}e^{\lambda t}\right)
+
\mathbf{K}
\left(\boldsymbol{\phi}e^{\lambda t}\right)
=\mathbf{0}.
\end{equation}

Colocando o fator exponencial em evidência,

\begin{equation}
\left(
\lambda^2\mathbf{M}
+
\lambda\mathbf{C}
+
\mathbf{K}
\right)
\boldsymbol{\phi}e^{\lambda t}
=\mathbf{0}.
\end{equation}

Como

\begin{equation}
e^{\lambda t}\neq 0
\end{equation}

para qualquer valor finito de $\lambda$ e $t$, podemos dividir a equação por
esse fator, resultando em

\begin{equation}
\boxed{
\left(
\lambda^2\mathbf{M}
+
\lambda\mathbf{C}
+
\mathbf{K}
\right)\boldsymbol{\phi}
=\mathbf{0}
}.
\label{eq:autovalores}
\end{equation}

Essa é a forma algébrica do problema de autovalores associado ao sistema.

\section{Condição para existência de solução não trivial}

A Equação~\eqref{eq:autovalores} é um sistema homogêneo. Uma solução sempre
possível é

\begin{equation}
\boldsymbol{\phi}=\mathbf{0},
\end{equation}

porém essa solução não representa um movimento físico do sistema.

Para existir uma solução não trivial,

\begin{equation}
\boldsymbol{\phi}\neq\mathbf{0},
\end{equation}

a matriz que multiplica $\boldsymbol{\phi}$ deve ser singular. Portanto,

\begin{equation}
\boxed{
\det\left(
\lambda^2\mathbf{M}
+
\lambda\mathbf{C}
+
\mathbf{K}
\right)=0
}.
\label{eq:determinante}
\end{equation}

A Equação~\eqref{eq:determinante} é a equação característica do sistema.

\section{Ordem do polinômio característico}

As matrizes $\mathbf{M}$, $\mathbf{C}$ e $\mathbf{K}$ possuem dimensão
$N\times N$. Cada elemento da matriz

\begin{equation}
\lambda^2\mathbf{M}
+\lambda\mathbf{C}
+\mathbf{K}
\end{equation}

é, no máximo, um polinômio de segundo grau em $\lambda$.

Ao calcular o determinante de uma matriz $N\times N$, o maior grau possível
resultante é

\begin{equation}
2N.
\end{equation}

Logo, o polinômio característico pode ser escrito como

\begin{equation}
\boxed{
P(\lambda)
=
a_{2N}\lambda^{2N}
+a_{2N-1}\lambda^{2N-1}
+\cdots
+a_1\lambda+a_0
=0
}.
\label{eq:polinomio-geral}
\end{equation}

Portanto, para um sistema com $N$ graus de liberdade, existem, em geral,
$2N$ raízes, contando multiplicidades e raízes complexas.

\section{Significado físico das raízes}

As raízes do polinômio característico são os valores de $\lambda$ que
determinam a resposta temporal do sistema.

Considere uma raiz complexa escrita como

\begin{equation}
\lambda=\alpha+i\omega_d,
\end{equation}

em que $\alpha$ é a parte real e $\omega_d$ é a parte imaginária.

A solução correspondente possui a forma

\begin{equation}
e^{\lambda t}
=
e^{(\alpha+i\omega_d)t}
=
e^{\alpha t}e^{i\omega_d t}.
\end{equation}

Assim, a parte real de $\lambda$ controla o comportamento da amplitude,
enquanto a parte imaginária está relacionada à oscilação.

\subsection*{Parte real negativa}

Se

\begin{equation}
\operatorname{Re}(\lambda)<0,
\end{equation}

então

\begin{equation}
e^{\operatorname{Re}(\lambda)t}\rightarrow 0
\qquad\text{quando}\qquad t\rightarrow\infty.
\end{equation}

Nesse caso, a resposta transitória decai com o tempo, caracterizando
comportamento amortecido e estável para o modo correspondente.

\subsection*{Parte real positiva}

Se

\begin{equation}
\operatorname{Re}(\lambda)>0,
\end{equation}

a amplitude cresce exponencialmente com o tempo. Isso indica instabilidade
dinâmica do modo correspondente.

\subsection*{Parte real nula}

Se

\begin{equation}
\operatorname{Re}(\lambda)=0,
\end{equation}

não existe decaimento exponencial da amplitude. Nesse caso, para uma raiz
puramente imaginária, o sistema apresenta comportamento oscilatório sem
amortecimento exponencial.

\subsection*{Parte imaginária}

Quando

\begin{equation}
\lambda=\alpha\pm i\omega_d,
\end{equation}

a presença de $\omega_d$ indica uma componente oscilatória. As raízes
complexas aparecem em pares conjugados quando as matrizes e os coeficientes do
problema são reais.

\section{Escolha de um sistema numérico com dois graus de liberdade}

Para cumprir as tarefas numéricas da APC2, será considerado um sistema com

\begin{equation}
N=2.
\end{equation}

Escolhemos os seguintes parâmetros físicos:

\begin{table}[H]
\centering
\begin{tabular}{ccc}
\toprule
Parâmetro & Símbolo & Valor \\
\midrule
Massa 1 & $m_1$ & $2\,\mathrm{kg}$ \\
Massa 2 & $m_2$ & $1\,\mathrm{kg}$ \\
Amortecimento 1 & $c_1$ & $2\,\mathrm{N\,s/m}$ \\
Amortecimento de acoplamento & $c_2$ & $1\,\mathrm{N\,s/m}$ \\
Amortecimento 2 & $c_3$ & $3\,\mathrm{N\,s/m}$ \\
Rigidez 1 & $k_1$ & $8\,\mathrm{N/m}$ \\
Rigidez de acoplamento & $k_2$ & $4\,\mathrm{N/m}$ \\
Rigidez 2 & $k_3$ & $5\,\mathrm{N/m}$ \\
\bottomrule
\end{tabular}
\caption{Parâmetros adotados para o sistema de dois graus de liberdade.}
\end{table}

Considerando a configuração usual de duas massas acopladas, as matrizes são

\begin{equation}
\mathbf{M}
=
\begin{bmatrix}
m_1 & 0\\
0 & m_2
\end{bmatrix}
=
\begin{bmatrix}
2 & 0\\
0 & 1
\end{bmatrix}.
\end{equation}

Para o amortecimento,

\begin{equation}
\mathbf{C}
=
\begin{bmatrix}
c_1+c_2 & -c_2\\
-c_2 & c_2+c_3
\end{bmatrix}.
\end{equation}

Substituindo os valores,

\begin{equation}
\mathbf{C}
=
\begin{bmatrix}
2+1 & -1\\
-1 & 1+3
\end{bmatrix}
=
\boxed{
\begin{bmatrix}
3 & -1\\
-1 & 4
\end{bmatrix}
}.
\end{equation}

Analogamente, para a rigidez,

\begin{equation}
\mathbf{K}
=
\begin{bmatrix}
k_1+k_2 & -k_2\\
-k_2 & k_2+k_3
\end{bmatrix}.
\end{equation}

Portanto,

\begin{equation}
\mathbf{K}
=
\begin{bmatrix}
8+4 & -4\\
-4 & 4+5
\end{bmatrix}
=
\boxed{
\begin{bmatrix}
12 & -4\\
-4 & 9
\end{bmatrix}
}.
\end{equation}

\section{Construção explícita do sistema de equações}

Para $N=2$,

\begin{equation}
\mathbf{x}(t)
=
\begin{bmatrix}
x_1(t)\\
x_2(t)
\end{bmatrix}.
\end{equation}

A equação matricial é

\begin{equation}
\begin{bmatrix}
2 & 0\\
0 & 1
\end{bmatrix}
\begin{bmatrix}
\ddot{x}_1\\
\ddot{x}_2
\end{bmatrix}
+
\begin{bmatrix}
3 & -1\\
-1 & 4
\end{bmatrix}
\begin{bmatrix}
\dot{x}_1\\
\dot{x}_2
\end{bmatrix}
+
\begin{bmatrix}
12 & -4\\
-4 & 9
\end{bmatrix}
\begin{bmatrix}
x_1\\
x_2
\end{bmatrix}
=
\begin{bmatrix}
0\\
0
\end{bmatrix}.
\end{equation}

Realizando as multiplicações matriciais, a primeira equação é

\begin{equation}
2\ddot{x}_1+3\dot{x}_1-\dot{x}_2+12x_1-4x_2=0.
\label{eq:eq1}
\end{equation}

A segunda equação é

\begin{equation}
\ddot{x}_2-\dot{x}_1+4\dot{x}_2-4x_1+9x_2=0.
\label{eq:eq2}
\end{equation}

Portanto, o sistema diferencial completo é

\begin{equation}
\boxed{
\begin{cases}
2\ddot{x}_1+3\dot{x}_1-\dot{x}_2+12x_1-4x_2=0,\\[2mm]
\ddot{x}_2-\dot{x}_1+4\dot{x}_2-4x_1+9x_2=0.
\end{cases}}
\end{equation}

\section{Obtenção do polinômio característico}

Utilizando

\begin{equation}
\mathbf{x}(t)=\boldsymbol{\phi}e^{\lambda t},
\qquad
\boldsymbol{\phi}
=
\begin{bmatrix}
\phi_1\\
\phi_2
\end{bmatrix},
\end{equation}

obtemos a matriz dinâmica

\begin{equation}
\mathbf{A}(\lambda)
=
\lambda^2\mathbf{M}
+\lambda\mathbf{C}
+\mathbf{K}.
\end{equation}

Substituindo as matrizes escolhidas,

\begin{equation}
\mathbf{A}(\lambda)
=
\lambda^2
\begin{bmatrix}
2 & 0\\
0 & 1
\end{bmatrix}
+
\lambda
\begin{bmatrix}
3 & -1\\
-1 & 4
\end{bmatrix}
+
\begin{bmatrix}
12 & -4\\
-4 & 9
\end{bmatrix}.
\end{equation}

Efetuando a soma termo a termo,

\begin{equation}
\mathbf{A}(\lambda)
=
\begin{bmatrix}
2\lambda^2+3\lambda+12
&
-\lambda-4\\
-\lambda-4
&
\lambda^2+4\lambda+9
\end{bmatrix}.
\end{equation}

A condição de existência de solução não trivial é

\begin{equation}
\det\left(\mathbf{A}(\lambda)\right)=0.
\end{equation}

Como a matriz possui dimensão $2\times2$,

\begin{equation}
\det
\begin{bmatrix}
a & b\\
c & d
\end{bmatrix}
=ad-bc.
\end{equation}

Logo,

\begin{align}
0
&=
(2\lambda^2+3\lambda+12)
(\lambda^2+4\lambda+9)
-
(-\lambda-4)(-\lambda-4).
\end{align}

Vamos desenvolver os produtos separadamente.

Primeiro,

\begin{align}
(2\lambda^2+3\lambda+12)
(\lambda^2+4\lambda+9)
&=
2\lambda^4
+8\lambda^3
+18\lambda^2\\
&\quad
+3\lambda^3
+12\lambda^2
+27\lambda\\
&\quad
+12\lambda^2
+48\lambda
+108.
\end{align}

Agrupando os termos semelhantes,

\begin{equation}
(2\lambda^2+3\lambda+12)
(\lambda^2+4\lambda+9)
=
2\lambda^4
+11\lambda^3
+42\lambda^2
+75\lambda
+108.
\end{equation}

Agora,

\begin{align}
(-\lambda-4)(-\lambda-4)
&=(\lambda+4)^2\\
&=\lambda^2+8\lambda+16.
\end{align}

Substituindo os dois resultados no determinante,

\begin{align}
0
&=
2\lambda^4
+11\lambda^3
+42\lambda^2
+75\lambda
+108\\
&\quad
-\left(\lambda^2+8\lambda+16\right).
\end{align}

Finalmente,

\begin{align}
0
&=
2\lambda^4
+11\lambda^3
+(42-1)\lambda^2
+(75-8)\lambda
+(108-16).
\end{align}

Assim,

\begin{equation}
\boxed{
2\lambda^4
+11\lambda^3
+41\lambda^2
+67\lambda
+92
=0
}.
\label{eq:polinomio-final}
\end{equation}

\section{Organização na forma padrão}

Comparando a Equação~\eqref{eq:polinomio-final} com

\begin{equation}
P(\lambda)
=
a_n\lambda^n
+a_{n-1}\lambda^{n-1}
+\cdots
+a_1\lambda
+a_0,
\end{equation}

temos

\begin{equation}
n=4
\end{equation}

e os coeficientes

\begin{align}
a_4 &= 2,\\
a_3 &= 11,\\
a_2 &= 41,\\
a_1 &= 67,\\
a_0 &= 92.
\end{align}

Portanto, o vetor de coeficientes que pode ser fornecido ao programa do
PPC2 é

\begin{equation}
\boxed{
\mathbf{a}
=
[2,\;11,\;41,\;67,\;92]
}.
\end{equation}

O polinômio na forma solicitada é

\begin{equation}
\boxed{
P(\lambda)
=
2\lambda^4
+11\lambda^3
+41\lambda^2
+67\lambda
+92
}.
\end{equation}

\section{Raízes do polinômio e interpretação física}

Para verificar o comportamento dinâmico do sistema, as raízes do polinômio
podem ser obtidas numericamente. Para o polinômio obtido, as raízes são
aproximadamente

\begin{align}
\lambda_1 &\approx -1{,}99793+2{,}40836i,\\
\lambda_2 &\approx -1{,}99793-2{,}40836i,\\
\lambda_3 &\approx -0{,}75207+2{,}03277i,\\
\lambda_4 &\approx -0{,}75207-2{,}03277i.
\end{align}

As raízes aparecem em pares complexos conjugados, como esperado para um
polinômio de coeficientes reais.

Todas as raízes possuem parte real negativa:

\begin{equation}
\operatorname{Re}(\lambda_i)<0.
\end{equation}

Consequentemente, os fatores exponenciais associados possuem decaimento com
o tempo. Os modos são, portanto, oscilatórios amortecidos.

As partes imaginárias,

\begin{equation}
2{,}40836
\qquad\text{e}\qquad
2{,}03277,
\end{equation}

representam as frequências angulares amortecidas dos dois pares de modos, em
$\mathrm{rad/s}$, enquanto as partes reais determinam a taxa de decaimento
exponencial das amplitudes.

É importante destacar que essa interpretação das raízes serve como conexão
entre a formulação matemática da APC2 e o comportamento físico do sistema,
enquanto a determinação das raízes pelo método de Bairstow será realizada na
etapa computacional do PPC2.

\section{Conclusão}

A partir da equação matricial do sistema massa--mola--amortecedor,

\begin{equation}
\mathbf{M}\ddot{\mathbf{x}}
+
\mathbf{C}\dot{\mathbf{x}}
+
\mathbf{K}\mathbf{x}
=0,
\end{equation}

foi adotada a solução exponencial

\begin{equation}
\mathbf{x}(t)=\boldsymbol{\phi}e^{\lambda t}.
\end{equation}

A substituição dessa solução levou ao problema algébrico

\begin{equation}
\left(
\lambda^2\mathbf{M}
+\lambda\mathbf{C}
+\mathbf{K}
\right)\boldsymbol{\phi}=0,
\end{equation}

e a condição para existência de solução não trivial resultou em

\begin{equation}
\det
\left(
\lambda^2\mathbf{M}
+\lambda\mathbf{C}
+\mathbf{K}
\right)=0.
\end{equation}

Para o sistema escolhido com dois graus de liberdade, foi obtido um polinômio
de quarta ordem,

\begin{equation}
\boxed{
P(\lambda)
=
2\lambda^4
+11\lambda^3
+41\lambda^2
+67\lambda
+92
}.
\end{equation}

Como $N=2$, verifica-se diretamente que a ordem do polinômio é

\begin{equation}
2N=4.
\end{equation}

O polinômio está organizado na forma padrão e seus coeficientes podem ser
utilizados diretamente na implementação do método de Bairstow no PPC2. As
raízes obtidas para o exemplo possuem partes reais negativas e partes
imaginárias não nulas, indicando modos oscilatórios amortecidos para os
parâmetros adotados.

\section*{Referências}

\begin{enumerate}
    \item GONTIJO, R. G. \textit{Notas de aula do curso de Cálculo Numérico
    Aplicado}. 2026.
    \item CHAPRA, S. C.; CANALE, R. P. \textit{Métodos Numéricos para
    Engenharia}. 5. ed. McGraw-Hill, 2008.
\end{enumerate}

\end{document}
